% Options for packages loaded elsewhere
% Options for packages loaded elsewhere
\PassOptionsToPackage{unicode}{hyperref}
\PassOptionsToPackage{hyphens}{url}
\PassOptionsToPackage{dvipsnames,svgnames,x11names}{xcolor}
%
\documentclass[
]{agujournal2019}
\usepackage{xcolor}
\usepackage{amsmath,amssymb}
\setcounter{secnumdepth}{5}
\usepackage{iftex}
\ifPDFTeX
  \usepackage[T1]{fontenc}
  \usepackage[utf8]{inputenc}
  \usepackage{textcomp} % provide euro and other symbols
\else % if luatex or xetex
  \usepackage{unicode-math} % this also loads fontspec
  \defaultfontfeatures{Scale=MatchLowercase}
  \defaultfontfeatures[\rmfamily]{Ligatures=TeX,Scale=1}
\fi
\usepackage{lmodern}
\ifPDFTeX\else
  % xetex/luatex font selection
\fi
% Use upquote if available, for straight quotes in verbatim environments
\IfFileExists{upquote.sty}{\usepackage{upquote}}{}
\IfFileExists{microtype.sty}{% use microtype if available
  \usepackage[]{microtype}
  \UseMicrotypeSet[protrusion]{basicmath} % disable protrusion for tt fonts
}{}
\makeatletter
\@ifundefined{KOMAClassName}{% if non-KOMA class
  \IfFileExists{parskip.sty}{%
    \usepackage{parskip}
  }{% else
    \setlength{\parindent}{0pt}
    \setlength{\parskip}{6pt plus 2pt minus 1pt}}
}{% if KOMA class
  \KOMAoptions{parskip=half}}
\makeatother
% Make \paragraph and \subparagraph free-standing
\makeatletter
\ifx\paragraph\undefined\else
  \let\oldparagraph\paragraph
  \renewcommand{\paragraph}{
    \@ifstar
      \xxxParagraphStar
      \xxxParagraphNoStar
  }
  \newcommand{\xxxParagraphStar}[1]{\oldparagraph*{#1}\mbox{}}
  \newcommand{\xxxParagraphNoStar}[1]{\oldparagraph{#1}\mbox{}}
\fi
\ifx\subparagraph\undefined\else
  \let\oldsubparagraph\subparagraph
  \renewcommand{\subparagraph}{
    \@ifstar
      \xxxSubParagraphStar
      \xxxSubParagraphNoStar
  }
  \newcommand{\xxxSubParagraphStar}[1]{\oldsubparagraph*{#1}\mbox{}}
  \newcommand{\xxxSubParagraphNoStar}[1]{\oldsubparagraph{#1}\mbox{}}
\fi
\makeatother


\usepackage{longtable,booktabs,array}
\usepackage{calc} % for calculating minipage widths
% Correct order of tables after \paragraph or \subparagraph
\usepackage{etoolbox}
\makeatletter
\patchcmd\longtable{\par}{\if@noskipsec\mbox{}\fi\par}{}{}
\makeatother
% Allow footnotes in longtable head/foot
\IfFileExists{footnotehyper.sty}{\usepackage{footnotehyper}}{\usepackage{footnote}}
\makesavenoteenv{longtable}
\usepackage{graphicx}
\makeatletter
\newsavebox\pandoc@box
\newcommand*\pandocbounded[1]{% scales image to fit in text height/width
  \sbox\pandoc@box{#1}%
  \Gscale@div\@tempa{\textheight}{\dimexpr\ht\pandoc@box+\dp\pandoc@box\relax}%
  \Gscale@div\@tempb{\linewidth}{\wd\pandoc@box}%
  \ifdim\@tempb\p@<\@tempa\p@\let\@tempa\@tempb\fi% select the smaller of both
  \ifdim\@tempa\p@<\p@\scalebox{\@tempa}{\usebox\pandoc@box}%
  \else\usebox{\pandoc@box}%
  \fi%
}
% Set default figure placement to htbp
\def\fps@figure{htbp}
\makeatother


% definitions for citeproc citations
\NewDocumentCommand\citeproctext{}{}
\NewDocumentCommand\citeproc{mm}{%
  \begingroup\def\citeproctext{#2}\cite{#1}\endgroup}
\makeatletter
 % allow citations to break across lines
 \let\@cite@ofmt\@firstofone
 % avoid brackets around text for \cite:
 \def\@biblabel#1{}
 \def\@cite#1#2{{#1\if@tempswa , #2\fi}}
\makeatother
\newlength{\cslhangindent}
\setlength{\cslhangindent}{1.5em}
\newlength{\csllabelwidth}
\setlength{\csllabelwidth}{3em}
\newenvironment{CSLReferences}[2] % #1 hanging-indent, #2 entry-spacing
 {\begin{list}{}{%
  \setlength{\itemindent}{0pt}
  \setlength{\leftmargin}{0pt}
  \setlength{\parsep}{0pt}
  % turn on hanging indent if param 1 is 1
  \ifodd #1
   \setlength{\leftmargin}{\cslhangindent}
   \setlength{\itemindent}{-1\cslhangindent}
  \fi
  % set entry spacing
  \setlength{\itemsep}{#2\baselineskip}}}
 {\end{list}}
\usepackage{calc}
\newcommand{\CSLBlock}[1]{\hfill\break\parbox[t]{\linewidth}{\strut\ignorespaces#1\strut}}
\newcommand{\CSLLeftMargin}[1]{\parbox[t]{\csllabelwidth}{\strut#1\strut}}
\newcommand{\CSLRightInline}[1]{\parbox[t]{\linewidth - \csllabelwidth}{\strut#1\strut}}
\newcommand{\CSLIndent}[1]{\hspace{\cslhangindent}#1}



\setlength{\emergencystretch}{3em} % prevent overfull lines

\providecommand{\tightlist}{%
  \setlength{\itemsep}{0pt}\setlength{\parskip}{0pt}}



 


\usepackage{url} %this package should fix any errors with URLs in refs.
\usepackage{lineno}
\usepackage[inline]{trackchanges} %for better track changes. finalnew option will compile document with changes incorporated.
\usepackage{soul}
\linenumbers
\makeatletter
\@ifpackageloaded{caption}{}{\usepackage{caption}}
\AtBeginDocument{%
\ifdefined\contentsname
  \renewcommand*\contentsname{Table of contents}
\else
  \newcommand\contentsname{Table of contents}
\fi
\ifdefined\listfigurename
  \renewcommand*\listfigurename{List of Figures}
\else
  \newcommand\listfigurename{List of Figures}
\fi
\ifdefined\listtablename
  \renewcommand*\listtablename{List of Tables}
\else
  \newcommand\listtablename{List of Tables}
\fi
\ifdefined\figurename
  \renewcommand*\figurename{Figure}
\else
  \newcommand\figurename{Figure}
\fi
\ifdefined\tablename
  \renewcommand*\tablename{Table}
\else
  \newcommand\tablename{Table}
\fi
}
\@ifpackageloaded{float}{}{\usepackage{float}}
\floatstyle{ruled}
\@ifundefined{c@chapter}{\newfloat{codelisting}{h}{lop}}{\newfloat{codelisting}{h}{lop}[chapter]}
\floatname{codelisting}{Listing}
\newcommand*\listoflistings{\listof{codelisting}{List of Listings}}
\makeatother
\makeatletter
\makeatother
\makeatletter
\@ifpackageloaded{caption}{}{\usepackage{caption}}
\@ifpackageloaded{subcaption}{}{\usepackage{subcaption}}
\makeatother
\usepackage{bookmark}
\IfFileExists{xurl.sty}{\usepackage{xurl}}{} % add URL line breaks if available
\urlstyle{same}
\hypersetup{
  pdftitle={Comparing Patient-Reported Outcomes Between Dupilumab and Conventional Therapy in Adults with Moderate-to-Severe Atopic Dermatitis},
  pdfauthor={Louie Neri},
  pdfkeywords={Atopic dermatitis, Dupilumab, Patient-reported
outcomes, Comparative effectiveness, Grant proposal},
  colorlinks=true,
  linkcolor={blue},
  filecolor={Maroon},
  citecolor={Blue},
  urlcolor={Blue},
  pdfcreator={LaTeX via pandoc}}


\journalname{NIH Grant Proposal}

\draftfalse

\begin{document}
\title{Comparing Patient-Reported Outcomes Between Dupilumab and
Conventional Therapy in Adults with Moderate-to-Severe Atopic
Dermatitis}

\authors{Louie Neri\affil{}}

\correspondingauthor{Louie Neri}{}


\begin{abstract}
This grant proposal compares patient-reported outcomes between dupilumab
and conventional therapy in adults with moderate-to-severe atopic
dermatitis using observational electronic health record data.
\end{abstract}

\section*{Plain Language Summary}
This project examines whether adults with moderate-to-severe atopic
dermatitis report better quality of life and fewer depressive symptoms
after starting dupilumab than after starting conventional therapy.




\section{Specific Aims}\label{specific-aims}

Atopic dermatitis (AD) is a chronic inflammatory skin condition
affecting an estimated 7\% to 10\% of adults in the United States
(Silverberg \& Hanifin, 2013). Moderate-to-severe AD produces persistent
pruritus, sleep disturbance, visible skin changes, and psychological
distress, which together reduce quality of life and elevate the risk of
depression and anxiety (Silverberg et al., 2018). The disease burden
experienced by patients extends well beyond what clinical severity
scores capture.

Dupilumab, a biologic therapy targeting IL-4 and IL-13 signaling, has
reshaped the treatment landscape for moderate-to-severe AD since its FDA
approval in 2017 (Simpson et al., 2016). The evidence supporting its
use, however, comes largely from randomized controlled trials with
restrictive eligibility criteria. These trials typically exclude older
adults, patients with multiple comorbidities, and those on concurrent
medications, which limits how well their findings translate to routine
clinical practice. Real-world data comparing patient-reported outcomes
(PROs) between dupilumab and conventional therapy remain sparse,
particularly over the 12-month time horizon relevant to clinical
decision-making (Afshari et al., 2024).

This gap matters for three reasons. First, clinicians and patients
making treatment decisions lack robust comparative data on outcomes that
reflect what patients actually experience. Second, cost differences
between biologic and conventional therapies are substantial.
Out-of-pocket costs for dupilumab can reach several thousand dollars per
year depending on insurance coverage, and the resulting financial burden
can independently worsen depressive symptoms even when skin symptoms
improve (Drucker et al., 2017). Third, understanding these trade-offs is
essential for shared decision-making that reflects both clinical
effectiveness and the patient-centered consequences of treatment choice.

The long-term goal of this research is to generate real-world evidence
on patient-centered outcomes that can inform treatment selection and
health policy for adults with moderate-to-severe AD. The objective of
this application is to compare changes in disease-specific quality of
life and depressive symptoms between patients initiating dupilumab
versus conventional therapy over 12 months, using electronic health
record (EHR) data from a large academic medical center. The underlying
rationale is that observational data reflecting real-world clinical
pathways complement trial evidence and identify outcomes meaningful to
patients, which supports the design of future clinical studies. This
proposal responds to PAR-24-280, Clinical Observational Studies in
Musculoskeletal, Rheumatic, and Skin Diseases, which seeks well-designed
observational research that improves understanding of disease symptoms
and outcomes important to patients with skin conditions.

The specific aims are:

\textbf{Aim 1: Compare 12-month changes in disease-specific quality of
life between dupilumab and conventional therapy.} Changes in the
Dermatology Life Quality Index (DLQI) (Finlay \& Khan, 1994) will be
evaluated between treatment groups, with confounding by indication
addressed through inverse probability of treatment weighting (IPTW). The
working hypothesis is that dupilumab-treated patients will show greater
DLQI improvement than those on conventional therapy, with differences
exceeding the minimal clinically important difference (MCID) of 4 points
(Basra et al., 2015).

\textbf{Aim 2: Compare 12-month changes in depressive symptoms between
dupilumab and conventional therapy.} Changes in the Patient Health
Questionnaire-9 (PHQ-9) (Kroenke et al., 2001) will be evaluated between
treatment groups using the same IPTW approach. The hypothesis is that
dupilumab-treated patients will show greater reductions in depressive
symptoms, consistent with the established relationship between skin
disease burden and psychological well-being.

\textbf{Aim 3: Examine whether treatment effects vary by age, sex, and
race.} Prespecified interaction terms will test for effect modification
by demographic characteristics. This aim will identify subgroups that
may experience differential benefits and will surface potential
disparities in treatment outcomes. Insurance status and socioeconomic
factors will be examined as exploratory moderators.

The expected output is real-world comparative effectiveness evidence on
patient-reported outcomes in AD that is directly relevant to clinical
decision-making. Secondary contributions include characterizing disease
symptom trajectories under routine care, supporting the use of DLQI and
PHQ-9 as endpoints in future clinical studies, and identifying sources
of treatment effect heterogeneity that can inform efforts to reduce
disparities in AD care. This work aligns with the NIAMS mission to
improve outcomes and function in patients with skin conditions.

\newpage{}

\section{Research Strategy}\label{research-strategy}

\subsection{Significance}\label{significance}

\subsubsection{Burden of atopic
dermatitis}\label{burden-of-atopic-dermatitis}

AD is the most common chronic inflammatory skin disease and affects a
substantial proportion of US adults (Silverberg \& Hanifin, 2013). The
condition is characterized by intense pruritus, erythematous and
excoriated lesions, and a relapsing-remitting course that meaningfully
diminishes quality of life. The associated sleep disturbance, social
stigma, impaired work productivity, and elevated rates of depression and
anxiety amplify this burden considerably (Yu et al., 2018). The
psychosocial impact of moderate-to-severe AD can be comparable to that
of asthma and diabetes (Eckert et al., 2017), and the outcomes patients
care about most are experiential: how the disease affects daily
activities, relationships, and emotional well-being rather than the
specific values on a clinical severity scale.

\subsubsection{Treatment options and the limits of current
evidence}\label{treatment-options-and-the-limits-of-current-evidence}

Conventional therapy includes topical corticosteroids, calcineurin
inhibitors, and systemic immunosuppressants such as cyclosporine,
methotrexate, and mycophenolate mofetil. These agents can reduce disease
severity but carry risks of long-term toxicity and yield variable
efficacy (Eichenfield et al., 2014). Dupilumab, the first FDA-approved
biologic for AD, has demonstrated efficacy in pivotal randomized trials
(Simpson et al., 2016). Those trials, however, enrolled highly selected
populations under controlled conditions. Their findings do not
necessarily extend to the broader population seen in routine practice,
where adherence is lower, comorbidities are common, and clinical
monitoring is less frequent.

\subsubsection{Gaps in patient-reported outcomes
evidence}\label{gaps-in-patient-reported-outcomes-evidence}

Despite growing adoption of dupilumab, head-to-head observational
comparisons of patient-reported outcomes between biologic and
conventional therapy in real-world settings remain limited. Most
existing studies rely on clinician-assessed severity scores such as EASI
and SCORAD (Chopra \& Silverberg, 2018). Clinical severity captures
important biological information but does not reflect the full scope of
what patients experience. DLQI was chosen for this study because it is a
condition-specific instrument that captures the particular impact of
skin disease on daily life, and because condition-specific measures tend
to be more responsive to treatment-related change than generic measures
when a real treatment effect is present. PHQ-9 adds a complementary
dimension by capturing depressive symptom burden, which is a recognized
comorbidity that amplifies overall disease impact. The pairing of these
two instruments produces a more complete picture of how treatments
affect patients than either clinical severity scores or generic health
measures alone.

\subsubsection{Financial toxicity and access
disparities}\label{financial-toxicity-and-access-disparities}

Out-of-pocket costs for dupilumab vary widely depending on insurance
coverage, ranging from minimal copayments to several thousand dollars
annually (Drucker et al., 2017). This variation creates selection
effects in who receives biologic therapy, with implications for both the
internal validity of observational comparisons and for health equity.
Financial burden can also independently worsen depression scores and
reduce adherence, potentially attenuating treatment benefits regardless
of pharmacologic efficacy. These cost-related dynamics must be accounted
for when interpreting PRO differences between treatment groups and when
designing future studies that address economic barriers to care.

\subsubsection{Importance to the field}\label{importance-to-the-field}

By generating the kind of observational evidence that PAR-24-280 is
designed to support, this project will provide data that directly inform
the design of future clinical studies in AD, including identification of
appropriate primary endpoints and characterization of the subgroups most
likely to benefit from biologic therapy.

\subsection{Innovation}\label{innovation}

This proposal advances the field in five respects.

\emph{Dual patient-reported outcome focus.} Existing comparative studies
of AD treatments have predominantly relied on clinician-assessed
severity endpoints (Chopra \& Silverberg, 2018). Centering this study on
two validated patient-reported instruments as co-primary outcomes shifts
the evidentiary focus toward what patients experience. DLQI captures
disease-specific impact on daily life and has strong responsiveness in
dermatology populations; PHQ-9 captures psychological burden and is
widely established for depression screening and severity assessment.
Pairing a condition-specific instrument with a mental health instrument
is uncommon in comparative AD research and reflects the multidimensional
nature of disease impact.

\emph{Integration of mental health outcomes in dermatology research.}
Depression is a recognized comorbidity of AD (Yu et al., 2018), but few
observational studies have systematically tracked depressive symptom
trajectories across treatment strategies. Including PHQ-9 as a primary
endpoint alongside DLQI makes the psychological dimension of treatment
benefit testable rather than assumed.

\emph{Systematic evaluation of treatment effect heterogeneity.} A single
average treatment effect can mask substantial variation across
subgroups. This study formally tests whether benefits differ by age,
sex, and race. Given documented disparities in AD treatment access
across demographic groups (Tackett et al., 2020), this approach
identifies subgroups that may experience differential benefits. Effect
modification will be assessed through formal interaction tests rather
than side-by-side comparison of significance labels across strata, the
latter being a common but invalid way to establish that effects truly
differ.

\emph{Methodological rigor in addressing confounding by indication.}
Patients receiving dupilumab typically have more severe, refractory
disease than those on conventional therapy. IPTW will create a
pseudo-population in which measured baseline covariates are balanced
across treatment groups, and covariate balance will be verified through
standardized mean differences (threshold: SMD \textless{} 0.1) and
propensity score overlap to evaluate positivity. E-value sensitivity
analyses will quantify how strong an unmeasured confounder would need to
be to explain away observed effects (VanderWeele \& Ding, 2017). The
combination provides a transparent framework for causal interpretation
in an observational setting.

\emph{Emphasis on clinical meaningfulness.} Interpretation will center
on effect size, confidence interval width, and compatibility with
established MCIDs (DLQI: 4 points, (Basra et al., 2015); PHQ-9: 5
points, (Kroenke et al., 2001)), rather than on whether a result crosses
a p-value threshold. Statistically significant findings do not guarantee
clinically important effects, particularly in large samples where small
differences can produce low p-values. Nonsignificant findings in smaller
samples may reflect imprecision rather than absence of effect.
Interpretation grounded in effect magnitude and clinical thresholds is
more useful for treatment decisions than a binary significance label.

\subsection{Approach}\label{approach}

\subsubsection{Overview and conceptual
model}\label{overview-and-conceptual-model}

This prospective cohort study will compare patient-reported outcomes
between adults with moderate-to-severe AD initiating dupilumab versus
conventional therapy. The conceptual model serves as the organizing
framework: treatment type (the primary exposure) affects disease
activity, itch intensity, and sleep disturbance, which in turn influence
DLQI and PHQ-9. Confounders including baseline disease severity,
treatment history, insurance status, comorbidities, and socioeconomic
factors affect both treatment selection and outcomes and must be
addressed through propensity score methods. The model assigns each
variable an explicit role as exposure, mediator, moderator, or
confounder, which prevents the analysis from devolving into a list of
variables without rationale and ensures every analytic decision ties
back to the hypothesized causal structure.

\subsubsection{Study design}\label{study-design}

A prospective cohort design is appropriate because randomizing patients
to dupilumab versus conventional therapy is neither practical nor
ethical. Dupilumab is typically prescribed after conventional options
have failed, and insurance authorization creates natural treatment
pathways that a research protocol cannot easily override. Two cohorts
will be enrolled at treatment initiation and followed for 12 months,
with PROs collected at baseline and at 3, 6, and 12 months. Prospective
collection ensures standardized outcome measurement rather than reliance
on inconsistently documented chart data.

The principal threat to validity is confounding by indication. Dupilumab
recipients tend to have more severe disease, meaning baseline
differences between groups could bias observed associations. This is a
structural problem that propensity score methods partially address for
measured confounders but cannot address for unmeasured ones, which is
why sensitivity analyses are essential.

\subsubsection{Study population and data
source}\label{study-population-and-data-source}

EHR data from a large academic medical center with an active dermatology
specialty clinic will be used. EHR data add clinical detail that
administrative claims often lack, including laboratory values, disease
severity markers, and clinical notes that are essential when outcome
ascertainment depends on timing and clinical context. The study
population includes adults (age 18 and older) with moderate-to-severe AD
confirmed by ICD-10 diagnostic codes and baseline EASI scores of 16 or
higher, who initiated either dupilumab or conventional systemic therapy
between 2020 and 2025. At least one baseline and one follow-up PRO
assessment are required. Exclusion criteria are concurrent use of other
biologic agents, participation in interventional clinical trials during
the study period, and diagnoses of other inflammatory skin conditions
that could confound outcome assessment.

DLQI and PHQ-9 will be extracted from routine clinical assessments
documented at baseline and follow-up visits. Both instruments have
established reliability, validity, and responsiveness, and both are
widely used in clinical practice, which helps maximize data availability
and minimize respondent burden. Additional data elements include
demographics (age, sex, race/ethnicity), EASI scores, treatment history
(number of prior systemic therapies), comorbidities including prior
depression/anxiety diagnoses, insurance type, and socioeconomic
indicators.

\subsubsection{Variables and
measurement}\label{variables-and-measurement}

\begin{longtable}[]{@{}
  >{\raggedright\arraybackslash}p{(\linewidth - 8\tabcolsep) * \real{0.2586}}
  >{\raggedright\arraybackslash}p{(\linewidth - 8\tabcolsep) * \real{0.1724}}
  >{\raggedright\arraybackslash}p{(\linewidth - 8\tabcolsep) * \real{0.2241}}
  >{\raggedright\arraybackslash}p{(\linewidth - 8\tabcolsep) * \real{0.2241}}
  >{\raggedright\arraybackslash}p{(\linewidth - 8\tabcolsep) * \real{0.1207}}@{}}
\caption{Study variables and
measurement}\label{tbl-variables}\tabularnewline
\toprule\noalign{}
\begin{minipage}[b]{\linewidth}\raggedright
Variable Type
\end{minipage} & \begin{minipage}[b]{\linewidth}\raggedright
Variable
\end{minipage} & \begin{minipage}[b]{\linewidth}\raggedright
Description
\end{minipage} & \begin{minipage}[b]{\linewidth}\raggedright
Measurement
\end{minipage} & \begin{minipage}[b]{\linewidth}\raggedright
Scale
\end{minipage} \\
\midrule\noalign{}
\endfirsthead
\toprule\noalign{}
\begin{minipage}[b]{\linewidth}\raggedright
Variable Type
\end{minipage} & \begin{minipage}[b]{\linewidth}\raggedright
Variable
\end{minipage} & \begin{minipage}[b]{\linewidth}\raggedright
Description
\end{minipage} & \begin{minipage}[b]{\linewidth}\raggedright
Measurement
\end{minipage} & \begin{minipage}[b]{\linewidth}\raggedright
Scale
\end{minipage} \\
\midrule\noalign{}
\endhead
\bottomrule\noalign{}
\endlastfoot
Primary Exposure & Treatment group & Dupilumab vs.~conventional therapy
& EHR prescription records & Categorical \\
Primary Outcome & Disease-specific QoL & DLQI change score (baseline to
12 months) & Range: 0-30 & Continuous \\
Primary Outcome & Depressive symptoms & PHQ-9 change score (baseline to
12 months) & Range: 0-27 & Continuous \\
Confounder & Age & Patient age at enrollment & Years & Continuous \\
Confounder & Sex & Male or female & Self-report & Categorical \\
Confounder & Race/ethnicity & Race/ethnic group & EHR demographics &
Categorical \\
Confounder & Baseline EASI score & Disease severity at enrollment &
Range: 0-72 & Continuous \\
Confounder & Treatment history & Prior systemic therapies & Count &
Continuous \\
Confounder & Disease duration & Years since AD diagnosis & Years &
Continuous \\
Confounder & Insurance status & Private, Medicare/Medicaid, or other &
EHR records & Categorical \\
Confounder & Depression/anxiety history & Documented prior diagnosis &
ICD codes & Categorical \\
\end{longtable}

\subsubsection{Mediating, moderating, and confounding
relationships}\label{mediating-moderating-and-confounding-relationships}

These three variable roles are frequently conflated but require
different analytic strategies. A mediator is caused by the exposure and
carries some portion of its effect to the outcome. Disease activity
(EASI score change), itch intensity, and sleep disturbance are the
hypothesized mediators in this framework. These will not be adjusted for
in the primary analysis because conditioning on a mediator blocks the
causal pathway through which treatment produces its effect, biasing the
estimate toward the null.

A moderator changes the strength or direction of the treatment-outcome
relationship and is sometimes called an effect modifier. Age, sex, and
baseline disease severity are the hypothesized moderators, examined
through prespecified interaction terms in Aim 3. A moderator does not
need to be associated with the exposure. It only needs to change the
exposure's effect on the outcome.

A confounder is associated with both treatment selection and outcomes
but does not lie on the causal pathway. Confounders must be measured
before exposure to avoid conditioning on post-treatment variables that
could introduce bias. Insurance status, disease duration, treatment
history, depression/anxiety history, and socioeconomic status are all
confounders and will be included in the propensity score model.

\subsubsection{Statistical analysis
plan}\label{statistical-analysis-plan}

\textbf{Propensity score estimation.} A propensity score is the
probability that a patient received a particular treatment given their
observed baseline characteristics. It compresses multiple confounders
into a single score between zero and one and functions as a balancing
score: patients with similar scores have, on average, similar
distributions of observed covariates regardless of which treatment they
received. Propensity scores will be estimated via logistic regression
with treatment group as the dependent variable and the following
baseline covariates as predictors: age, sex, race/ethnicity, baseline
EASI score, number of prior systemic therapy failures, disease duration,
insurance type, depression/anxiety history, and relevant comorbidities.
Variable selection is based on subject matter knowledge of likely
confounders. After weighting, covariate balance will be assessed using
standardized mean differences (threshold: SMD \textless{} 0.1) and
propensity score overlap will be inspected visually to evaluate the
positivity assumption, which requires a nonzero probability of receiving
each treatment for every combination of covariates (Austin, 2011).

\textbf{Primary analysis.} The average treatment effect on the treated
(ATT) will be estimated using IPTW in weighted linear regression models.
IPTW creates a pseudo-population in which measured covariates are
balanced across treatment groups, approximating what randomization would
achieve. Primary outcomes are 12-month change scores for DLQI and PHQ-9.
Point estimates will be reported with 95\% confidence intervals, and
clinical significance will be evaluated against established MCIDs (DLQI:
4 points; PHQ-9: 5 points), focusing on whether confidence intervals
exclude these thresholds and what effect sizes remain compatible with
the data. Hypothesis testing at alpha = 0.05 will be reported, though a
point estimate paired with its confidence interval is more informative
for guiding treatment decisions than a binary significance label.

\textbf{Effect modification analysis (Aim 3).} Prespecified interaction
terms for age, sex, and race will be included in regression models.
Effect modification is established through the interaction test itself,
not by comparing significance labels across subgroups. Stratum-specific
effect sizes and confidence intervals will be reported to characterize
the pattern of heterogeneity. Insurance status and socioeconomic factors
will be examined as exploratory moderators.

\textbf{Missing data.} Missing data will be handled using multiple
imputation by chained equations (MICE), under the assumption that data
are missing at random conditional on observed covariates (White et al.,
2011). Because EHR data often exhibit missingness that is not completely
at random, sensitivity analyses will compare imputed results to
complete-case analyses. Substantial divergence between approaches will
prompt more cautious interpretation.

\textbf{Sensitivity and robustness analyses.} Three robustness checks
are planned: (1) propensity score matching (1:1 nearest-neighbor with
caliper of 0.2 standard deviations of the logit of the propensity score)
as an alternative to weighting; (2) doubly robust estimation combining
IPTW with outcome regression, which remains consistent as long as at
least one of the two models is correctly specified; and (3) E-value
calculations to quantify the minimum strength of association that an
unmeasured confounder would need to have with both treatment and outcome
to explain away observed effects (VanderWeele \& Ding, 2017). The three
assumptions underlying causal inference (exchangeability, positivity,
and consistency) will be explicitly evaluated and discussed.

\subsubsection{Influence of age, sex, and
race}\label{influence-of-age-sex-and-race}

Demographic characteristics are likely to shape both treatment patterns
and outcomes. Older versus younger adults may have different baseline
quality of life and depressive symptom profiles, and age-related
variation in insurance coverage and formulary restrictions may affect
the probability of receiving dupilumab. Sex differences may emerge in
how patients perceive and report quality of life and mood symptoms over
follow-up. Race may relate to disparities in treatment access, with
differences in insurance coverage, referral patterns, and prior
authorization requirements creating selection effects in who receives
biologic versus conventional therapy (Tackett et al., 2020). All three
variables will enter the propensity score model as covariates, and
interaction analyses (Aim 3) will evaluate whether treatment effects
differ across demographic subgroups. Stratum-specific results will be
interpreted through effect sizes and confidence intervals rather than
significance labels alone.

\subsubsection{Potential challenges and
solutions}\label{potential-challenges-and-solutions}

\begin{longtable}[]{@{}
  >{\raggedright\arraybackslash}p{(\linewidth - 2\tabcolsep) * \real{0.3667}}
  >{\raggedright\arraybackslash}p{(\linewidth - 2\tabcolsep) * \real{0.6333}}@{}}
\caption{Potential challenges and proposed
solutions}\label{tbl-challenges}\tabularnewline
\toprule\noalign{}
\begin{minipage}[b]{\linewidth}\raggedright
Challenge
\end{minipage} & \begin{minipage}[b]{\linewidth}\raggedright
Proposed Solution
\end{minipage} \\
\midrule\noalign{}
\endfirsthead
\toprule\noalign{}
\begin{minipage}[b]{\linewidth}\raggedright
Challenge
\end{minipage} & \begin{minipage}[b]{\linewidth}\raggedright
Proposed Solution
\end{minipage} \\
\midrule\noalign{}
\endhead
\bottomrule\noalign{}
\endlastfoot
Confounding by indication: dupilumab patients have more severe,
refractory disease & IPTW with thorough covariate balance assessment
(SMD \textless{} 0.1); doubly robust estimation as sensitivity analysis;
E-value calculations for unmeasured confounding \\
Missing PRO data from incomplete follow-up or inconsistent EHR
documentation & Multiple imputation by chained equations with
sensitivity analyses comparing complete-case results; evaluation of
whether missingness patterns relate to patient characteristics or
outcome risk \\
Selection bias from differential out-of-pocket costs affecting who
receives dupilumab & Include insurance status and socioeconomic
indicators in propensity score model; assess residual imbalance
post-weighting \\
Treatment switching or discontinuation during follow-up &
Intention-to-treat analysis as primary approach; per-protocol analysis
as sensitivity check with documented switching rates and reasons \\
Insufficient sample size for subgroup analyses by race & Report
stratum-specific estimates with confidence intervals; acknowledge
precision limitations transparently; avoid interpreting nonsignificance
as absence of effect \\
\end{longtable}

\subsubsection{Study timeline}\label{study-timeline}

\begin{longtable}[]{@{}lcccc@{}}
\caption{Study timeline}\label{tbl-timeline}\tabularnewline
\toprule\noalign{}
Activity & Year 1 Q1-Q2 & Year 1 Q3-Q4 & Year 2 Q1-Q2 & Year 2 Q3-Q4 \\
\midrule\noalign{}
\endfirsthead
\toprule\noalign{}
Activity & Year 1 Q1-Q2 & Year 1 Q3-Q4 & Year 2 Q1-Q2 & Year 2 Q3-Q4 \\
\midrule\noalign{}
\endhead
\bottomrule\noalign{}
\endlastfoot
IRB approval and data access agreements & ✓ & & & \\
Cohort identification and data extraction & ✓ & ✓ & & \\
Propensity score estimation and balance assessment & & ✓ & & \\
Primary analyses (Aims 1 and 2) & & ✓ & ✓ & \\
Effect modification analyses (Aim 3) & & & ✓ & \\
Sensitivity analyses and robustness checks & & & ✓ & \\
Manuscript preparation and dissemination & & & ✓ & ✓ \\
\end{longtable}

\newpage{}

\section{References}\label{references}

\phantomsection\label{refs}
\begin{CSLReferences}{1}{0}
\vspace{1em}

\bibitem[\citeproctext]{ref-afshari_unraveling_2024}
Afshari, M., Kolackova, M., Rosecka, M., Čelakovská, J., \& Krejsek, J.
(2024). Unraveling the skin; a comprehensive review of atopic
dermatitis, current understanding, and approaches. \emph{Front Immunol},
\emph{15}, 1361005. \url{https://doi.org/10.3389/fimmu.2024.1361005}

\bibitem[\citeproctext]{ref-austin_introduction_2011}
Austin, P. C. (2011). An {Introduction} to {Propensity} {Score}
{Methods} for {Reducing} the {Effects} of {Confounding} in
{Observational} {Studies}. \emph{Multivariate Behav Res}, \emph{46}(3),
399--424. \url{https://doi.org/10.1080/00273171.2011.568786}

\bibitem[\citeproctext]{ref-basra_determining_2015}
Basra, M. K. A., Salek, M. S., Camilleri, L., Sturkey, R., \& Finlay, A.
Y. (2015). Determining the minimal clinically important difference and
responsiveness of the {Dermatology} {Life} {Quality} {Index} ({DLQI}):
Further data. \emph{Dermatology}, \emph{230}(1), 27--33.
\url{https://doi.org/10.1159/000365390}

\bibitem[\citeproctext]{ref-chopra_assessing_2018}
Chopra, R., \& Silverberg, J. I. (2018). Assessing the severity of
atopic dermatitis in clinical trials and practice. \emph{Clin Dermatol},
\emph{36}(5), 606--615.
\url{https://doi.org/10.1016/j.clindermatol.2018.05.012}

\bibitem[\citeproctext]{ref-drucker_burden_2017}
Drucker, A. M., Wang, A. R., Li, W.-Q., Sevetson, E., Block, J. K., \&
Qureshi, A. A. (2017). The {Burden} of {Atopic} {Dermatitis}: {Summary}
of a {Report} for the {National} {Eczema} {Association}. \emph{J Invest
Dermatol}, \emph{137}(1), 26--30.
\url{https://doi.org/10.1016/j.jid.2016.07.012}

\bibitem[\citeproctext]{ref-eckert_impact_2017}
Eckert, L., Gupta, S., Amand, C., Gadkari, A., Mahajan, P., \& Gelfand,
J. M. (2017). Impact of atopic dermatitis on health-related quality of
life and productivity in adults in the {United} {States}: {An} analysis
using the {National} {Health} and {Wellness} {Survey}. \emph{J Am Acad
Dermatol}, \emph{77}(2), 274--279.e3.
\url{https://doi.org/10.1016/j.jaad.2017.04.019}

\bibitem[\citeproctext]{ref-eichenfield_guidelines_2014}
Eichenfield, L. F., Tom, W. L., Berger, T. G., Krol, A., Paller, A. S.,
Schwarzenberger, K., et al. (2014). Guidelines of care for the
management of atopic dermatitis: Section 2. {Management} and treatment
of atopic dermatitis with topical therapies. \emph{J Am Acad Dermatol},
\emph{71}(1), 116--132. \url{https://doi.org/10.1016/j.jaad.2014.03.023}

\bibitem[\citeproctext]{ref-finlay_dermatology_1994}
Finlay, A. Y., \& Khan, G. K. (1994). Dermatology {Life} {Quality}
{Index} ({DLQI})--a simple practical measure for routine clinical use.
\emph{Clin Exp Dermatol}, \emph{19}(3), 210--216.
\url{https://doi.org/10.1111/j.1365-2230.1994.tb01167.x}

\bibitem[\citeproctext]{ref-kroenke_phq-9_2001}
Kroenke, K., Spitzer, R. L., \& Williams, J. B. W. (2001). The {PHQ}-9.
\emph{J GEN INTERN MED}, \emph{16}(9), 606--613.
\url{https://doi.org/10.1046/j.1525-1497.2001.016009606.x}

\bibitem[\citeproctext]{ref-silverberg_adult_2013}
Silverberg, J. I., \& Hanifin, J. M. (2013). Adult eczema prevalence and
associations with asthma and other health and demographic factors: A
{US} population-based study. \emph{J Allergy Clin Immunol},
\emph{132}(5), 1132--1138.
\url{https://doi.org/10.1016/j.jaci.2013.08.031}

\bibitem[\citeproctext]{ref-silverberg_patient_2018}
Silverberg, J. I., Gelfand, J. M., Margolis, D. J., Boguniewicz, M.,
Fonacier, L., Grayson, M. H., et al. (2018). Patient burden and quality
of life in atopic dermatitis in {US} adults: {A} population-based
cross-sectional study. \emph{Ann Allergy Asthma Immunol}, \emph{121}(3),
340--347. \url{https://doi.org/10.1016/j.anai.2018.07.006}

\bibitem[\citeproctext]{ref-simpson_two_2016}
Simpson, E. L., Bieber, T., Guttman-Yassky, E., Beck, L. A., Blauvelt,
A., Cork, M. J., et al. (2016). Two {Phase} 3 {Trials} of {Dupilumab}
versus {Placebo} in {Atopic} {Dermatitis}. \emph{N Engl J Med},
\emph{375}(24), 2335--2348. \url{https://doi.org/10.1056/NEJMoa1610020}

\bibitem[\citeproctext]{ref-tackett_structural_2020}
Tackett, K. J., Jenkins, F., Morrell, D. S., McShane, D. B., \&
Burkhart, C. N. (2020). Structural racism and its influence on the
severity of atopic dermatitis in {African} {American} children.
\emph{Pediatr Dermatol}, \emph{37}(1), 142--146.
\url{https://doi.org/10.1111/pde.14058}

\bibitem[\citeproctext]{ref-vanderweele_sensitivity_2017}
VanderWeele, T. J., \& Ding, P. (2017). Sensitivity {Analysis} in
{Observational} {Research}: {Introducing} the {E}-{Value}. \emph{Ann
Intern Med}, \emph{167}(4), 268--274.
\url{https://doi.org/10.7326/M16-2607}

\bibitem[\citeproctext]{ref-white_multiple_2011}
White, I. R., Royston, P., \& Wood, A. M. (2011). Multiple imputation
using chained equations: Issues and guidance for practice. \emph{Stat
Med}, \emph{30}(4), 377--399. \url{https://doi.org/10.1002/sim.4067}

\bibitem[\citeproctext]{ref-yu_systematic_2018}
Yu, S. H., Drucker, A. M., Lebwohl, M., \& Silverberg, J. I. (2018). A
systematic review of the safety and efficacy of systemic corticosteroids
in atopic dermatitis. \emph{J Am Acad Dermatol}, \emph{78}(4),
733--740.e11. \url{https://doi.org/10.1016/j.jaad.2017.09.074}

\end{CSLReferences}




\end{document}
